\section{Discussion and conclusion}

The experiments support a specific claim: measurement-accessible tangent geometry can be used as a physically grounded preconditioner, and in several regimes it is more useful than the full state-space metric or a Euclidean gradient. They do not support the stronger claim that AQNG is a universally superior optimizer. Random-rank AQNG is often close to aligned AQNG, Adam remains competitive in the large/deep comparison, and the best readout orientation depends on architecture, depth, and system size.

The distinction between \emph{accessibility} and \emph{trainability} is central. A readout can capture nearly all of the relevant tangent covariance while the supervised gradient is effectively absent, as in the large-$n$ $U(1)$ control. Conversely, a readout can retain more tangent mass while producing a less useful preconditioner because the metric becomes poorly conditioned, as in the weight-two ablation. Retention is therefore a resource descriptor, not an optimization objective by itself. A more complete trainability picture requires at least three ingredients: accessible orientation, stable metric spectrum, and surviving task-gradient signal.

This also clarifies the relation to barren plateaus~\cite{McClean2018}. AQNG rescales and rotates an observed gradient; it cannot create supervised information that has already fallen below the estimation noise floor. In an exact arithmetic model, an inverse metric can compensate some common scaling between gradient and geometry, but finite damping, small-eigenvalue instability, and metric-estimation noise limit that mechanism. Demonstrating barren-plateau mitigation would require a dedicated scaling study of gradient variance or update signal-to-noise with substantially larger initialization ensembles and unclipped update diagnostics. The present $U(1)$ experiment is therefore a gradient-collapse control rather than a formal barren-plateau theorem.

Several resource caveats are also important. Cross-fitted alignment requires calibration tangents, and those measurements are part of the experimental cost. The finite-shot full-QNG control uses an analytic metric oracle, so it is not a hardware-cost-matched comparator. The primary 20-seed benchmark keeps the dataset split fixed, meaning that it probes initialization and optimization stochasticity rather than split-to-split generalization uncertainty. Finally, the exact projector identity used by AQNG assumes a jointly measurable commuting readout; extending the method to noncommuting observable families requires an explicit measurement-grouping model.

The readout-order ablation suggests a concrete next algorithmic step. Rather than choosing the subspace that maximizes retained tangent mass, one can optimize a regularized criterion that trades accessibility against spectral conditioning, for example through eigenvalue truncation, adaptive damping, or a condition-number penalty. A second direction is semantic structure: if the retained observables are tied to known input features or physical quantities, accessible eigendirections may be annotated with domain meaning rather than only geometric meaning. A third is a dedicated trainability-scaling experiment in which the raw gradient, accessible spectrum, metric-estimation noise, and shot complexity are tracked simultaneously.

AQNG converts measurement-accessible tangent geometry from a diagnostic into an optimizer. Its value is not universal dominance over first-order methods; it is the explicit coupling between the preconditioning geometry and the measurement interface. The experiments show that accessibility, conditioning, and task-gradient signal must be assessed separately. When all three are favorable, the accessible metric can provide competitive or improved optimization; when either conditioning or gradient signal fails, high retained tangent mass alone is insufficient.
